\section{Discussion}

\subsection{The Cumulative Arc}

The six experiments form a coherent chain of pre-registered
rejections that drive the experimental design. E1 lag-1 fails at
K=200 and the natural smoothing fix (E2) is falsified. E3
introduces change-point gating and avoids the K=200 hazard at
moderate K=50 cost; E4 shows no single magnitude threshold can
reach the joint feasibility region. E5 reaches feasibility via a
per-K vector but does so by approximating the static gate, with no
K=200 improvement. E6 introduces CUSUM, breaks the static-gate
collapse at K=200 with a $+0.9$~pp gain, and exposes the
single-parameter K=50 over-firing as the new failure mode.

This is not a series of separate findings; it is a single
research question --- ``what is the deployable recalibration
recipe?'' --- evaluated through a pre-registered sequence in which
each next protocol's choice is justified by the previous protocol's
rejection.

\subsection{Why Smoothing Fails (E2)}

The naive prior was that averaging the lag-1 calibration target
across multiple past windows would reduce variance and recover
some K=200 lag-1 benefit. The data falsifies the prior cleanly:
EWMA-3 produces $\rho = -0.94$ at K=200, worse than lag-1's
$-0.66$. The mechanism is bias-variance trade-off under
distributional shift. At high capacity a single window remediates
$\sim 8\%$ of the unpatched backlog and preferentially the
high-EPSS pairs, so the residual distribution at window $w$ is
materially different from $w - 3, w - 2, w - 1$. Averaging
introduces bias from older states; the bias dominates the variance
reduction. The result rules out simple history-based smoothing in
the studied regime.

\subsection{Why a Single $\tau$ Cannot Reach Feasibility (E3, E4)}

The single-magnitude detector at one $\tau$ value has a
one-parameter family of fire patterns. E3 at $\tau = 0.05$
satisfies the K=200 hazard tolerance (H1.E3) but fires too eagerly
at K=50 (H2.E3 cost $-5.3$~pp). E4 sweeps $\tau$ and shows the
Pareto curve sits entirely outside the joint feasibility region
defined by Paper~5/E3's two tolerances. The K=50 cost is uniformly
$\geq 4$~pp.

The structural reason: at the studied bi-weekly cadence, the
per-window $|\Delta_w|$ values are bimodal --- either "large"
(typically $> 0.10$) or "small" (typically $< 0.05$), with little
mass between. A single threshold therefore behaves approximately
binarily, and the K=50 cost is approximately constant across all
$\tau$ values that map to the same "mostly-fire" or
"mostly-don't-fire" regime.

\subsection{The Capacity-Aware Collapse (E5)}

The per-K threshold vector $\tau_K = \{0.20, 0.05, 0.02\}$ reaches
the joint feasibility region: the large $\tau_{50}$ never fires,
the small $\tau_{200}$ fires every window. The two extremes
collapse to lag-1 and fixed respectively, matching the static gate
exactly. At K=100 the detector intervenes 40\% of the time and
loses $0.4$~pp on average to lag-1 --- within bootstrap noise.

The substantive interpretation is that the magnitude detector
family cannot improve on the static gate without a richer
firing rule, because matching the static gate's K=200 decisions
requires firing at every K=200 window, and matching its K=50
decisions requires never firing at K=50.

\subsection{The CUSUM Win (E6)}

CUSUM accumulates evidence across windows and refrains from firing
on single-window noise excursions. At K=200, the $w = 4$ window has
$|\Delta_4| = 0.020 < k = 0.04$, so the accumulator stays at $0$
and CUSUM uses the lag-1 fit at W4. The lag-1 fit at K=200 W4 is
beneficial (the same $+3.4$~pp benefit that adaptive05 captured by
not firing in E3); CUSUM captures it without firing. The W5 and
intermediate windows similarly avoid the gate's revert-to-fixed
decision when cumulative evidence is below $h$.

The $+0.9$~pp K=200 gain over the static gate is the only
operationally meaningful improvement in the six-experiment sequence.
The single-parameter K=50 over-firing failure (H1.E6) is the
natural next problem; a per-K $(k_K, h_K)$ CUSUM vector under
pre-registration is the obvious follow-up paper.

\subsection{Operational Recipe}

Sharpened across all six experiments:
\begin{itemize}
\item \textbf{$K \leq 100$:} deploy lag-1 online recalibration
directly (E1, $\rho \approx 1.0$). No detector, no gate. The
calibration arc reduces to the simple ``refit on the previous
window'' recipe at these capacities.
\item \textbf{$K = 200$:} deploy CUSUM($k = 0.04, h = 0.10$) (E6,
$+0.9$~pp over the static gate). This is the only operationally
meaningful gain over the static gate in the program.
\item \textbf{Mixed-capacity deployments:} use a per-K CUSUM
$(k_K, h_K)$ vector (untested in this paper, identified as Future
work).
\end{itemize}

\subsection{The Honest Negatives Are the Contribution}

The pre-registered rejections in Table~\ref{tab:hypothesis_ledger}
are not failures of the work; they are the work. Each rejection
constrains the next experiment's design choice. The cumulative
finding is that the deployable-recalibration question has a
\textit{capacity-conditional} answer: lag-1 below the gate
threshold, CUSUM above it, no universal procedure across the
studied grid. Reporting this honestly under pre-registration
discipline is what distinguishes the work from a tuning exercise.

\subsection{What This Closes and Opens}

\textit{Closes:} the single-detector deployability question for the
magnitude and CUSUM families. The K=200 hazard problem identified
in E1 is solved by CUSUM (E6) at a $+0.9$~pp margin over the static
gate; the smoothing alternative (E2) is falsified; the
single-$\tau$ family is provably infeasible (E4); the capacity-aware
$\tau$ family collapses to the gate (E5).

\textit{Opens:} (i) per-K CUSUM $(k_K, h_K)$ vector under
pre-registration --- the direct H1.E6 fix; (ii) two-sided CUSUM,
EWMA-CUSUM, Bayesian online change-point detectors; (iii) real
fleet validation, particularly to test whether the noise-excursion
patterns CUSUM exploits in the synthetic generator generalise;
(iv) closed-loop calibration where the selection policy is also
recalibrated, addressing the selection-policy coupling threat that
all six experiments inherit.
